\ifx\danslelivre\undefined
\documentclass[10pt]{../fichiers-configuration-latex/smfart}
\usepackage{ifthen}
\newboolean{english}
\newboolean{francais}
\setboolean{english}{false}
\setboolean{francais}{true}
\input{../fichiers-configuration-latex/paquets}
\selectlanguage{francais}
\input{../fichiers-configuration-latex/lemme-proposition-theoreme}
\input{../fichiers-configuration-latex/numerotation}
\input{../fichiers-configuration-latex/macros}


\externaldocument{../13-theoreme-de-finitude/finitude}
\externaldocument{../16-classes-de-Chern/gysin}

\begin{document}

\author{Yves Laszlo}

\title{XIX. Un contre-exemple}
\date{exposés oraux des  11 et 18 juin 2008}

\setcounter{tocdepth}{3}

\begin{abstract}L'exposé est destiné à construire, suivant Gabber, un exemple d'immersion ouverte
$j:\ U\rightarrow X$ de schémas noethériens telle que $R^1j_*\ZZ/2\ZZ$ ne soit pas constructible.
Ceci donne un contre-exemple à une version du théorème de constructibilité de Gabber
(\refext{XIII}{13-constructibilité}) sans hypothèse de quasi-excellence. Cette construction donne
aussi un exemple de diviseur $D$, dont $U$ est le complémentaire en l'occurrence, dans une surface
régulière localement à croisements normaux (au sens de de Jong) mais pas globalement (\ref{19-dcn}).
\end{abstract}

\maketitle

\tableofcontents

\else
\setboolean{english}{false}
\setboolean{francais}{true}
\selectlanguage{francais}
\chapter{Exemples et contre-exemples}
\addtocontents{toc}{par Yves Laszlo}
\begin{center}
Yves Laszlo
\end{center}
\fi

\section{Introduction}

L'exposé est destiné à construire, suivant Gabber, un exemple d'immersion ouverte $j:\ U\rightarrow
X$ de schémas noethériens telle que $R^1j_*\ZZ/2\ZZ$ ne soit pas constructible. Ceci montre que
l'hypothèse de quasi-excellence du théorème de constructibilité de Gabber
(\refext{XIII}{13-constructibilité}) est indispensable. D'un point de vue géométrique, la
construction est intéressante : $U$ est le complémentaire d'un diviseur $D$ dans une surface
régulière $X$  mais possède une infinité de points doubles ordinaires ; en particulier, son lieu
régulier n'est pas ouvert ce qui lui interdit d'être quasi-excellent. Ce diviseur est un exemple de
diviseur dans une surface régulière localement à croisements normaux (au sens de de Jong) mais pas
globalement (\ref{19-dcn}).

\section{La construction}\label{19-cex}

On va construire une immersion ouverte $j:U\hookrightarrow X$ de schémas noethériens tel que
$R^1j_*(\ZZ/2\ZZ)$ ne soit pas constructible. En particulier, $X$ n'est pas quasi-excellent d'après
le théorème de constructibilité de Gabber (\refext{XIII}{13-constructibilité}). Cette dernière
propriété est d'ailleurs évidente par construction puisque $X$ contient un fermé intègre de
dimension $1$ ayant une infinités de points doubles ! Suivant \textbf{SGA 4}, on note ici
$k\{\underline{x}\}$ l'hensélisé à l'origine d'un anneau de polynômes $k[\underline{x}]$. On choisit
un corps parfait $k$ de caractéristique différente de $2$, au plus dénombrable, qui a la propriété
suivante : toute extension finie de $k$ admet un élément qui n'est pas un carré. On note $\bar k$
une clôture algébrique de $k$.  Par exemple, on peut prendre pour $k$ un corps fini ou un corps de
nombres. Dans la suite, on note $\Lambda=\ZZ/2\ZZ$.




 On commence par regarder le plan $\mathbf{A}^2=\Spec(k[x,y])$ privé des points génériques des
courbes  irréductibles ne coupant pas la droite $\Delta=\Spec(k[x])$ d'équation $y=0$. Ces courbes
sont exactement les courbes irréductibles d'équation $u(1+yg(x,y)),u\in k^*$. On pose donc
$$A_0=(1+yk[x,y])^{-1}k[x,y].$$ Le morphisme de localisation $k[x,y]\ra A_0$ identifie $\Spec(A_0)$
au sous-ensemble du plan $\mathbf{A}^2=\Spec(k[x,y])$ cherché. Les points de $\Spec(A_0)$ sont de
trois sortes \begin{itemize} \item Le point générique de $\mathbf{A}^2$ ; \item Les points
génériques des courbes irréductibles du plan qui rencontrent $\Delta$ ; \item Les points de
$\Spec(A_0)$ fermés dans $\mathbf{A}^2$ (qui sont les points fermés de $\Delta$ comme on va le
voir).  \end{itemize}


Notons qu'un point générique d'une courbe $C$ qui coupe $\Delta$ se spécialise dans $\Spec(A_0)$ sur
n'importe quel point fermé de $C\cap \Delta$ et donc n'est pas fermé dans $\Spec(A_0)$.

Par ailleurs, un point de $\Specmax(A_0)$ est donc défini par $(\bar x,\bar y)\in\bar k$. Si $\bar
y$ est non nul, étant algébrique sur $k$, son inverse est dans $k[\bar y]$ ce qui entraîne $(\bar
x,\bar y)\not\in\Spec(A_0)(\bar k)$.  L'immersion fermée
$$\Spec(k[x])=\Spec(A_0/yA_0)\hookrightarrow\Spec(A_0)$$ induit donc un homéomorphisme
$\Specmax(k[x])\iso\Specmax(A_0)$.


Si $\xi\in\Specmax(A_0)$, on note $\pi_\xi\in k[x]$ le générateur unitaire des polynômes nuls en
$\xi$ et on choisit une racine de $\pi_\xi$ dans $\bar k$ définissant un point géométrique $\bar\xi$
au dessus de $\xi$. On le voit comme un élément de $A_0$ via le plongement tautologique
$k[x]\hookrightarrow A_0$. Le couple $(\pi_\xi,y)$ est un système de coordonnées locales de $A_0$ en
$\xi$, \emph{i.e.} on a un isomorphisme\footnote{\label{19-abus-coord}On devrait plutôt dire que le
morphisme $k[X,Y]_{(0,0)}\ra A_{0,\xi}$ qui envoie $X$ sur $\pi_\xi$ et $Y$ sur $y$ induit un unique
isomorphisme $k(\xi)\{\pi_\xi,y\}\iso A_{0,\xi}^{\h}$.} \begin{equation}\label{19-isom-coord}
k(\xi)\{\pi_\xi,y\}\iso A_{0,\xi}^{\h}.  \end{equation}


Comme $k$ est dénombrable, $\Specmax(A_0)$ est infini dénombrable comme $\Specmax(k[x])$. Soit
$[i],i\geq 0$ la suite des idéaux maximaux de $A_0$.



\begin{lemme}\label{19-lemme-la-constr} Il existe une tour d'extensions quadratiques $$A_n\subset
A_{n+1}=A_n[\sqrt{f_n}],\ n\geq 0$$ et une suite de points $\xi_n\in\Specmax(A_n)$ d'images
$\zeta_n=\xi_n\cap A_0$ dans $\Spec(A_0)$ telle que \begin{itemize} \item Pour tout $m\leq n$ et
tout $\m\in\Spec(A_n)$ au dessus de $[m]$, l'image de $f_n$ dans $k(\m)$ n'est pas un carré.  \item
Pour tout $m< n$ et tout $\m\in\Spec(A_n)$ au dessus de $\xi_m$, l'image de $f_n$ dans $k(\m)$ n'est
pas un carré.  \item L'extension $A_n/A_0$ est étale en $\xi_n$ et $$f_n\equiv\pi_n^2+y\mod\xi_n^3$$
où $\pi_n=\pi_{\zeta_n}$.  \end{itemize} \end{lemme}



\begin{proof} Soit $n\geq 0$. Supposons construits $A_m,f_m,\xi_m,\zeta_m,0\leq m<n$ vérifiant les
propriétés précédentes et construisons $f_n\in A_n$ et $\xi_n\in\Spec(A_n)$. Si $n>0$, on pose
$A_n=A_{n-1}[\sqrt{f_{n-1}}]$ (si $n=0$ l'anneau $A_n$ est déjà construit).

Observons tout de suite que, la caractéristique de $k$ étant différente de $2$, l'extension
$A_n/A_0$ est génériquement étale.  Comme elle est finie et que $\Specmax(A_0)=\Specmax(k[x])$, le
lieu $L_1$ des points de $\Specmax(A_n)$ où $A_n/A_0$ n'est pas étale est de cardinal fini et de
complémentaire $L_1^c$ infini. De même, l'ensemble des points $L_2$ de $\Spec(A_n)$ au dessus des
$[m],m\leq n$ et des $\zeta_m,m<n$ est fini.

Choisissons alors $\xi_n\in L_1^c-L_2$ et notons $\zeta_n=\xi_n\cap A_0$ son image dans
$\Spec(A_0)$. Par construction, on a $$\zeta_n\not\in\{[m],m\leq n,\zeta_m,m<n\}\text{ et
}A_n/A_0\text{ est étale en }\xi_n.$$

Soit $\m\in L_2$. Comme $k(\m)$ est une extension finie de $k$ (théorème des zéros), on peut choisir
$a_\m\in k(\m)$ qui n'est pas un carré. Comme $\xi_n\not\in L_2$, les idéaux de $L_2\cup\{\xi_n^3\}$
sont deux à deux étrangers. Le lemme chinois assure qu'il existe $f_n\in A$ tel que $$f_n\equiv
a_\m\mod\m\text{ pour tout }\m\in L_2\text{ et }f_n\equiv \pi_n^2+y\mod\xi_n^3.$$

De tels $f_n,\xi_n,\zeta_n$ conviennent.\end{proof}

On pose $A=\limind\ A_n$. Soit $n\in\NN$. Comme l'extension d'anneaux intègres $A_n\hookrightarrow
A$ est entière, on peut choisir un point géométrique $\overline{\xi_{n,\infty}}$ de $\Spec(A)$ au
dessus de $\xi_n$. Il définit donc des points géométriques $\overline{\xi_n}$ (resp.
$\overline{\zeta_n}$) au dessus de $\xi_n$ (resp. $\zeta_n$).

Par construction, l'inclusion $A_{n+1}\hookrightarrow A$ définit un isomorphisme (cf. la
note~\ref{19-abus-coord}) $$\bar k\{\pi_n,y,z\}/(z^2-f_n)\iso A^{\hs}_{\overline{\xi_{n,\infty}}}$$
compatible à (\ref{19-isom-coord}), \emph{i.e.} tel que le diagramme $$\xymatrix{\bar
k\{\pi_n,y,z\}/(z^2-f_n)\ar[r]^-{\thicksim}& A^{\hs}_{\overline{\xi_{n,\infty}}}\\ \bar
k\{\pi_n,y\}\ar[r]^-{\thicksim}\ar[u]& A_{0,\overline{\zeta_n}}^{\hs}\ar[u] }$$ commute (rappelons
que l'hensélisation stricte commute aux limites inductives filtrantes, cf. \ega{IV}{4}{18.8.18}).

Avec ces préparatifs, on peut énoncer le résultat principal.

\begin{proposition}\label{19-prop-cexgabber} Soit $j$  l'immersion ouverte
$\Spec(A[1/y])\hookrightarrow \Spec(A)$ et $\eta$ le point générique de $D=V(y)$.  \begin{itemize}
\item $A$ est noethérien.  \item La dimension (sur $\FF_2$) de
$(R^1j_*\Lambda)_{\overline{\xi_{n,\infty}}}$ est $2$, alors que la dimension de
$(R^1j_*\Lambda)_{\bar\eta}$  est $1$.  \item En particulier, $R^1j_*\Lambda$ n'est pas
constructible.  \end{itemize} \end{proposition}

\begin{remarque}\label{19-rem-cexgabber}Notons que le diviseur (intègre)  $D=V(y)$ de la surface
régulière $\Spec(A)$ admet chaque $\xi_{n,\infty}$ comme point double (ordinaire). Il n'est donc pas
quasi-excellent puisque son lieu régulier (ou normal, c'est la même chose ici) n'est pas ouvert.
Quitte à éventuellement localiser $A$, on obtient alors un contre-exemple à la constructibilité avec
un schéma ambiant régulier (mais certes pas excellent) ! \end{remarque}

Le point \textbf{iii} découle immédiatement des points 1) et 2). Passons à la preuve des deux
premiers points.  Le reste de l'exposé est destiné à prouver les points \textbf{i} et \textbf{ii},
seuls points restant à montrer.

\section{Noethérianité de $A$} On va adapter (cf. proposition~\ref{19-crit-noeth}) à la situation
(en l'utilisant)  le critère usuel de noethérianité des limites inductives  que l'on rappelle :

\begin{theoreme}[\egazero{III}{10.3.1.3}]\label{19-ega-noeth}Soit $(A_i,\m_i)$ un système inductif
filtrant d'anneaux locaux noethériens. On suppose que tous les $A_i$ noethériens et que les
morphismes de transitions sont locaux et plats. Alors, si pour tout $i\leq j$, on a $\m_iA_j=\m_i$,
alors $\limind A_i$ est noethérien.\end{theoreme}

On utilisera sans le rappeler ensuite le critère de noethérianité de Cohen (\cite[3.4]{Nag62}) :

\begin{proposition}[Cohen]\label{19-copremier} Un anneau est noethérien si et seulement si tout
idéal premier est de type fini.\end{proposition}


Soit $A_i,i\geq 0$ un système inductif d'anneaux et $A_\infty=\limind A_i$. On suppose que les
morphismes $A_i\ra A_{i+1}$ sont finis et injectifs et que chaque $A_i$ est noethérien (ou, ce qui
revient au même, que $A_0$ est noethérien). En particulier, $\Spec(A_{i+1})\ra\Spec(A_i)$ est fini
et surjectif et  $\Spec(A_\infty)\ra\Spec(A_i)$ est entier et surjectif pour tout $i$ ce qu'on
utilisera sans plus de précaution. Leurs  fibres sont de dimension nulle.  Pour $\p\in\Spec(A_0)$,
on note $\tilde\star_\p$ la propriété

\textbf{Propriété $\tilde\star_\p$:}\emph{ Il existe $i$ tel que pour tout $j\geq i$ et tout
$\q\in\Spec(A_j)$ au dessus de $\p$, l'idéal $\q A_\infty$ est premier.}

\begin{proposition}\label{19-crit-noeth0} $A_\infty$ est noethérien si et seulement si  $A_0$ est
noethérien et tout idéal premier $\p$ de $A_0$ vérifie la
propriété $\tilde\star_\p$.
\end{proposition}

\begin{proof}On note $f:\Spec(A_\infty)\ra\Spec(A_0)$. Prouvons le premier point.

\underline{Suffisance}. Soit $\q_\infty\in\Spec(A_\infty)$ et $\p$ son image dans $\Spec(A_0)$.
Montrons que $\q_\infty$ est de type fini. Choisissons $i$ comme dans $\tilde\star_\p$ et soit
$\q=\q_\infty\cap A_i$. On a d'une part $\q A_\infty\subset\q_\infty$ et, d'autre part
$$\q_\infty\cap A_i\subset\q A_\infty\subset\q_\infty$$ ce qui assure l'égalité $\p=\q_\infty\cap
A_0=\q A_\infty\cap A_0$ de sorte que $\q A_\infty$ se spécialise sur  $\q_\infty$ dans $f^{-1}(\p)$
qui est de dimension $0$. On a donc $\q_\infty=\q A_\infty$ ce qui prouve que $\q_\infty$ est de
type fini comme $\q$ et on invoque~\ref{19-copremier}.

\underline{Nécessité}. Supposons $A_\infty$ noethérien et soit $\p\in\Spec(A_0)$. La fibre
$$f^{-1}(\p)=\Spec(A_\infty\otimes_{A_0}\kappa(\p))$$ est noethérienne de dimension nulle, donc de
cardinal fini. Comme $A_\infty$ est noethérien, on peut donc supposer que tous les idéaux premiers
de $f^{-1}(\p)$ sont engendrés par des éléments de $A_i$ pour $i$ convenable. Soit alors
$\q\in\Spec(A_j), j\geq i$ au dessus de $\p$.  Soit $\q'\in\Spec(A_\infty)$ au dessus de $\q$. Comme
$\q'$ est engendré par $\q'\cap A_i$, il l'est par $\q'\cap A_j=\q$, de sorte $\q A_\infty=\q'$ qui
est donc premier.\end{proof}



\begin{proposition}\label{19-crit-noeth} On garde les hypothèses et les notations de
\ref{19-crit-noeth0}. Si de plus les extensions $A_{i+1}/A_{i}$ sont \emph{plates}, $A_\infty$ est
noethérien si et seulement si tout idéal \emph{maximal} $\m$ de $A_0$ vérifie la propriété
$\tilde\star_\m$.\end{proposition}

\begin{proof}La nécessité découle du premier point.  D'après~\ref{19-crit-noeth0}, il suffit de
prouver la suffisance. Soit donc $\p\in\Spec(A_0)$ et montrons que $\p$ vérifie $\star_\p$.


\begin{lemme}Sous les conditions de la proposition, la propriété $\star_\p$ est équivalente à la
\textbf{Propriété $\star_\p$ : } \emph{Il existe $i$ tel que pour tout $l\geq j\geq i$ et tout
$\q\in\Spec(A_j)$ au dessus de $\p$, l'idéal $\q A_{l}$ est premier.} \end{lemme} \begin{proof}
Supposons $\star_\p$ vérifiée. Soit alors $\q\in\Spec(A_j)$ au dessus de $\p$.  On déduit déjà
$1\not\in\q A_\infty$. De plus, si $xy\in\q A_\infty$, il existe $l\geq j$, tel que $x,y\in A_l$.
Donc, $xy\in\q A_\infty\cap A_l=\q A_l$ (fidèle platitude de $A_\infty/A_l$) et donc par exemple
$x\in\q A_l\subset\q A_\infty$. On a donc $\star_\p\Rightarrow\tilde\star_\p$. L'autre implication
est tautologique.\end{proof}

La clef est de constater que la condition $\star_\p$ ne dépend que des fibres schématiques de $f_i :
\Spec(A_i)\ra\Spec(A_0)$ et donc est \emph{invariante par changement de base}, ce qui va permettre
de se ramener au cas local pour appliquer (\ref{19-ega-noeth}). Précisons.

\begin{lemme}\label{19-star-fib} Soit $\p\in\Spec(A_0)$. Les deux propriétés suivantes sont
équivalentes.  \begin{itemize} \item  La propriété $\star_\p$ est satisfaite.  \item Il existe $i$
tel que pour tout $l\geq j\geq i$ le morphisme induit $$\phi : f_l^{-1}(\p)\ra f_{j}^{-1}(\p)$$
entre les fibres schématiques soit bijectif à fibres réduites.  \end{itemize} \end{lemme}
\begin{proof} Supposons  $\star_\p$  vérifiée et choisissons $i\leq j\leq l$ comme dans $\star_\p$.
Soit $\q\in f_j^{-1}(\p)$ et posons $A=A_j/\q$ et $B=A_l/\q A_l$.  La fibre schématique
$\phi^{-1}(\q)$ est le $k(\q)=\Frac(A)$-schéma $\Spec(B\otimes_A \Frac(A))$. Comme $\q A_l$ est
premier, $B$ est intègre et donc $B\otimes_A \Frac(A)$ également (la tensorisation par $\Frac(A)$
est une localisation). Comme $A_l/A_j$ est finie, l'extension $B/A$ est finie de sorte que
$B\otimes_A \Frac(A)$ est à la fois de dimension finie sur $\Frac(A)$ et intègre, donc c'est un
corps, ce qui entraîne la seconde condition.

Inversement, supposons que $\phi^{-1}(\q)$ soit le spectre d'un corps. Autrement dit,  on suppose
que $B\otimes_A\text{Frac}(A)$ est intègre et on veut montrer que $B$ est intègre. Mais on a
\begin{sslemme} Soient $\phi : A\ra B$ un morphisme plat d'anneaux noethériens. Supposons $A$
intègre. Alors, $B$ est intègre si et seulement si la fibre générique $B\otimes_A\Frac(A)$ est
intègre.\end{sslemme} \begin{proof} Seul la suffisance demande explication.  Comme $\phi$ est plat,
il envoie idéaux associés (resp. minimaux) sur l'unique point associé de $\Spec(A)$, son point
générique ($A$ est intègre). Comme la fibre générique $B\otimes_A\text{Frac}(A)$ est un localisé,
les points associés de la fibre générique s'identifient par localisation aux points associés de
$\Spec(B)$. Comme cette fibre générique  est intègre, $\Spec(B)$ n'a qu'un point générique $\eta$
(ce qui assure l'irréductibilité de $\Spec(B)$), dont l'anneau local $B_\eta$ est un localisé de
$B\otimes_A\text{Frac}(A)$. Ainsi $B$ est sans point immergé et génériquement réduit ($S_1$ et
$R_0$) donc réduit d'où finalement intègre.  \end{proof} \end{proof} \medskip

Choisissons $\m$ maximal dans $A_0$ contenant $\p$ et choisissons $i$ comme dans $\star_\m$. Soit
$\m_i\in\Specmax(A_i)$ au dessus de $\m$ (($\Spec(A_i)\ra\Spec(A_0)$ fini et surjectif) et relevons
$\p\subset\m$ en $\p_i\subset\m_i$ dans $A_i$ (platitude).

Par construction, l'idéal $\m_iA_j$ est premier pour tout $j\geq i$. Mais on a $\m_iA_j\cap
A_i=\m_i$ ($A_j/A_i$ est fidèlement plate) de sorte que $\m_iA_j$ est maximal ($A_j/A_i$ est finie).
Soit alors $A'=(A_i)_{\m_i}$ et changeons de base par $A_0\ra A'$. La fibre spéciale du morphisme
fini $$\Spec(A'_j)\ra\Spec(A'_i),j\geq i$$ est le spectre du \emph{corps} $A_j/\m_iA_j=A_j/\m_j$ de
sorte que $A'_j,i\geq j$ est \emph{local} et que $A'_i\ra A'_j$ est local. Comme les fibres de
$\Spec(A'_j)\ra \Spec(A'_0)$ sont des fibres de $\Spec(A_j)\ra\Spec(A_0)$, le système inductif
$(A'_j)$ vérifie $\star_{\m'}$ pour tout $\m'\in\Specmax(A'_0)$ (\ref{19-star-fib}) de sorte que
$\limind A'_j$ est noethérien (\ref{19-ega-noeth}). D'après la proposition~\ref{19-crit-noeth0}, la
condition $\star_{\p'}$ est vérifiée pour tout $\p'\in\Spec(A'_0)$. Comme la condition $\star$ ne
dépend que des fibres (\ref{19-star-fib}), on déduit que $\star_{\q}$ est vérifiée pour tout $\q$
dans l'image de $\Spec(A'_i)\ra\Spec(A_0)$. Or, par construction $\p$ est dans l'image, ce qui
termine la preuve du point \textbf{i} de la proposition \ref{19-prop-cexgabber}.\end{proof} \medskip



\section{Étude des points doubles} Reste à prouver le point \textbf{ii} de la
proposition~\ref{19-prop-cexgabber}.

Comme au voisinage de $\eta$ l'immersion fermée $D=V(y)\hookrightarrow\Spec(A)$ est une immersion
d'un diviseur régulier dans une schéma régulier, le théorème de pureté
(\refext{XVI}{16-corollaire-purete-absolue-dcn}) assure que la dimension de
$(R^1j_*\Lambda)_{\bar\eta}$ est $1$.

Pour alléger les notations, on pose $x=\pi_n, f=f_n$ et $R=\bar k\{x,y,z\}/(z^2-f)$ et on rappelle
la congruence $$f=x^2+y+O(3)$$ où $O(n)$ désigne un polynôme de degré $\geq n$. On doit démontrer
l'énoncé suivant

\begin{proposition} La dimension de $H^1(\Spec(R[1/y],\Lambda)$ est $2$.\end{proposition}

D'après le théorème de pureté,\refext{XVI}{16-corollaire-purete-absolue-dcn}, la proposition est
conséquence du lemme géométrique bien connu suivant

\begin{lemme}[Lemme de Morse]\label{19-Morse} Il existe un système de coordonnée $X,y,Z$  de $\bar
k\{x,y,z\}$ tel que  $z^2-f=Z^2-X^2-y$.\end{lemme}


\begin{proof} On a $$z^2-f=z^2-x^2-y+O(3).$$ Posons $x_1=\sqrt{-1}x,x_2$. Dans ces coordonnées, on a
$$F:=z^2-f+y=\sum_{i,j}g_{i,j}x_ix_j$$ où $(g_{i,j})\in M_2(R)$ est une matrice symétrique de taille
2 à coefficients dans $R$. On a $$g=(g_{i,j})=\Id+O(1).$$ Considérons alors
$M=\mathbf{A}^{\frac{2(2+1)}{2}}$ paramétrant les matrices carrées symétriques. La différentielle de
$$\mu:\ \left\{ \begin{array}{ccc} M & \rightarrow& M \\ P &\mapsto & {}^tPP \\ \end{array}
\right.$$  en $P=\Id$ est le double de l'application identique, donc $\mu$ est étale en ce point. Le
point $g\in M(R)$ se relève donc en $P\in M(R)$. Autrement dit, on peut écrire $g={}^tPP$.  Le
morphisme ${}^t(x_i)\mapsto P{}^t(x_i)={}^t(y,X_1,Z_1)$ est étale et définit donc des coordonnées de
$\bar k\{x,y,z\}$. Dans ces coordonnées, on a $$Z^2-f+y=X_1^2+Z_1^2$$ et il suffit de poser
$X=\sqrt{-1}X_1, Z=Z_1$ pour conclure.\end{proof}

\medskip Rappelons que la caractéristique de $k$ est différente de $2$.  Ceci termine la preuve de
la proposition~\ref{19-prop-cexgabber}.

\section{$D$ est localement mais pas globalement un diviseur à croisements
normaux}\label{19-sec-dcn}

Commençons par une définition. Dans cette section $D$ désigne un diviseur effectif d'un schéma
régulier $X$ et $j:\ U=X-D\hookrightarrow X$ l'immersion ouverte du complémentaire.

\begin{definition} On conserve les notations précédentes.  \begin{itemize} \item On dit que $D$ est
localement un diviseur à croisements normaux (en abrégé, \emph{localement dcn}) si pour tout $x\in
D$, le localisé de Zariski $\Spec(\OO_{D,x})$ est un diviseur à croisements normaux de
$\Spec(\OO_{X,x})$.  \item Supposons $D$ localement dcn. On note $\varepsilon(x)$ le nombre de
branches analytiques de $\Spec(\OO_{D,x})$ et $\zeta(x)$ son nombre de composantes irréductibles. La
fonction $\varepsilon:\ x\mapsto\varepsilon(x)$ (resp. $\zeta:\ x\mapsto\zeta(x)$) est appelée
fonction de comptage analytique (resp. fonction de comptage Zariski).  \end{itemize}
\end{definition}

Avec les notations précédentes, si $\bar x$ est un point géométrique au dessus de $x\in D$ avec $D$
localement à croisements normaux, l'hensélisé strict $D_{(\bar x)}$ est un diviseur à croisements
normaux stricts de  $D_{(\bar x)}$. On a alors la caractérisation suivante :

\begin{lemme}\label{19-carac-fonction-dcn}  Avec les notations précédentes, supposons de plus que
$D$ est localement dcn et $\Lambda=\ZZ/\ell\ZZ$ avec $\ell$ un nombre premier inversible sur $X$.
Alors, les propositions suivantes sont équivalentes.  \begin{itemize} \item $R^1j_*\Lambda$ est
constructibles ; \item $Rj_*\Lambda$ est constructible ; \item la fonction de comptage analytique
$\varepsilon$ est constructible.  \end{itemize} \end{lemme}

\begin{proof}D'après le théorème de pureté (\refext{XVI}{16-corollaire-purete-absolue-dcn}), la
fibre $(Rj_*\Lambda)_{\bar x}$ est l'algèbre extérieure sur $$(R^1j_*\Lambda)_{\bar
x}=\Lambda^{\varepsilon(x)}.$$ Le lemme en découle immédiatement grâce à la caractérisation des
faisceaux constructibles à fibres finies (\sga{4}{IX}{prop. 2.13 (iii)}).\end{proof}

L'intérêt de ce lemme réside dans la proposition  suivante.

\begin{proposition} Avec les notations précédentes, supposons de plus que $D$ est localement dcn.
Alors, $\varepsilon$ est constructible si et seulement si $D$ est un diviseur à croisement
normaux.\end{proposition}

\begin{proof} Grâce au théorème de pureté (\refext{XVI}{16-corollaire-purete-absolue-dcn}) et au
lemme précédent, il suffit de prouver la partie directe. Supposons donc $\varepsilon$ constructible
et montrons que $D$ est à croisements normaux. Soit $\bar x$ un point géométrique au dessus de $x\in
D$. Puisque $D_{(\bar x)}$ est un diviseur à croisements normaux stricts, il existe un voisinage
étale $\pi:\ X'\rightarrow X$ de $x$ dans $X$, tel que le diviseur $D'=\pi^{-1}(D)$ est la somme de
diviseurs $D'_i$ qui sont réguliers en $x$ et qui se coupent transversalement en $x'$, image de
$\bar x$ dans $X'$. La fonction de comptage analytique $\varepsilon'$ de $D'$ est la somme des
fonctions de comptage analytiques $\varepsilon'_i$. Comme $\varepsilon'$   ne dépend que de
l'hensélisé strict, on a donc $$\varepsilon'=\varepsilon\circ\pi=\sum\varepsilon'_i.$$ En
particulier, $\varepsilon'$ est constructible comme $\varepsilon$. La fonction de comptage Zariski
$\zeta'$ de $D'$ certainement constructible de sorte que la différence $\varepsilon'-\zeta'$ l'est
aussi. Par hypothèse, $\varepsilon'-\zeta'$ s'annule sur $\spec{\OO_{X',x'}}$, donc sur l'ensemble
des générisations de $x'$. Comme elle est constructible, elle est nulle sur un voisinage ouvert $U'$
(Zariski) de $x'$. Comme $\epsilon'_i\geq\zeta_i'$, on a $\epsilon'_i=\zeta_i'$ sur $U'$ de sorte
que, quitte à restreindre $U'$, chaque diviseur $D_i$ est régulier sur $U'$. En se restreignant au
localisé strict de chaque point de $U'$, sur lequel on sait que $D'$ est un diviseur à croisements
normaux, on obtient que les $D_i$ se coupent transversalement de sorte que  la restriction de $D'$ à
$U'$ est un diviseur à croisements normaux stricts.\end{proof}
 
\begin{remarque} L'argument précédent appliqué à $\zeta$ assure que si le localisé Zariski de $D$ en
tout point est un diviseur à croisements normaux strict alors $D$ est un diviseur à croisement
normaux strict.\end{remarque}

Avec les notations de la proposition~\ref{19-prop-cexgabber} , on a donc obtenu le résultat suivant.

\begin{corollaire}\label{19-dcn} Le diviseur $D$ de la surface régulière $\spec{A}$ est localement à
croisements normaux mais pas globalement.\end{corollaire}

\ifx\danslelivre\undefined
\bibliography{../fichiers-configuration-latex/bibliographie}
\bibliographystyle{../fichiers-configuration-latex/smfalpha}
\end{document}
\fi
